% ----------------------------------------------------------------------------
%                                  INTRODUÇÃO
% ----------------------------------------------------------------------------

\chapter{INTRODUÇÃO}
\label{chap:introducao}

O Transtorno de Déficit de Atenção com Hiperatividade (TDAH) é um transtorno do neurodesenvolvimento caracterizado por níveis prejudiciais de desatenção, desorganização e/ou hiperatividade-impulsividade \cite{apa2014}. Frequentemente, o TDAH apresenta comorbidade com o Transtorno Opositor Desafiador (TOD), marcado por um padrão persistente de humor raivoso, comportamento desafiador e índole vingativa. Segundo \citeonline{barkley2015}, o manejo clínico e pedagógico dessas condições exige intervenções multimodais, uma vez que as disfunções executivas e de regulação emocional inerentes a esses transtornos impactam a rotina familiar e as relações sociais da criança.

No contexto educacional brasileiro, o direito ao acompanhamento adequado para o público neurodivergente é respaldado de forma contundente pela legislação. De acordo com a lei nº 14.254 \citeonline{brasil2021lei14254},as escolas devem oferecer acompanhamento integral para educandos com TDAH e outros transtornos, exigindo a articulação direta entre as escolas e as famílias. Segundo \citeonline{brasil2015lei13146}, a Lei Brasileira de Inclusão da Pessoa com Deficiência Lei nº 13.146/2015, evidencia-se a exigência institucional por práticas pedagógicas que garantam o pleno desenvolvimento cognitivo e social do aluno.

Apesar do amparo legal, a prática cotidiana nas escolas esbarra na dificuldade de estabelecer uma rede de apoio eficiente. A pesquisa parte da constatação de que a falta de comunicação eficaz entre família e escola compromete severamente o desempenho acadêmico e social dessas crianças, além de gerar conflitos e dificultar a aplicação de intervenções consistentes. Autores da psicologia educacional \cite{dessen2007}; \cite{epstein2011} corroboram que a parceria família-escola atua como um fator de proteção essencial, sendo que a ausência de um alinhamento tático entre pais e professores rotineiramente agrava os episódios de oposição do aluno.

Nesse cenário de descompasso, as Tecnologias de Informação e Comunicação (TICs) emergem como ferramentas viabilizadoras da educação inclusiva. O sistema proposto busca oferecer funcionalidades para o registro de comportamentos, o acompanhamento de estratégias pedagógicas e o compartilhamento de informações relevantes, mitigando ruídos na comunicação \cite{almeida2018}. Contudo, a maioria das plataformas educacionais disponíveis atualmente carece de um embasamento focado em saúde mental ou impõe barreiras financeiras incompatíveis com a realidade brasileira.

Diante das lacunas tecnológicas observadas, este trabalho propõe o desenvolvimento de um sistema web voltado ao auxílio da comunicação integrada para o manejo comportamental de crianças com TDAH e TOD. O projeto é conduzido e aplicado no âmbito do próprio Instituto Federal de Brasília (IFB), instituição que servirá como ambiente de pesquisa para a avaliação e validação empírica. Os resultados esperados incluem a melhoria da comunicação entre os atores envolvidos, a redução de conflitos e o fortalecimento dos vínculos sociais e acadêmicos dos alunos, favorecendo um processo educativo efetivamente mais inclusivo e colaborativo.

\vspace{0.3cm}
\section{Tema}

O tema deste trabalho consiste no desenvolvimento de um sistema de informação voltado ao auxílio no manejo comportamental de crianças diagnosticadas com Transtorno de Déficit de Atenção com Hiperatividade (TDAH) e Transtorno Opositor Desafiador (TOD), com foco em promover a integração comunicativa entre pais e professores.

\vspace{0.3cm}
\section{Problema}

Diante da falta de ferramentas acessíveis que integrem a rede de apoio escolar e familiar, e considerando as limitações metodológicas dos aplicativos atuais, a questão que norteia este projeto é: Como um sistema web pode apoiar estratégias de manejo comportamental de crianças com TDAH e TOD no contexto educacional e familiar?

\vspace{0.3cm}
\subsection{Objetivo Geral}

O objetivo geral deste trabalho é desenvolver um sistema de informação que auxilie no manejo comportamental de crianças diagnosticadas com TDAH e TOD, promovendo a integração comunicativa eficaz entre pais e professores, de modo a melhorar o acompanhamento e a intervenção no comportamento dessas crianças.

\vspace{0.3cm}
\subsection{Objetivos Específicos}

\begin{itemize}
    \item Contextualizar as características e impactos do TDAH e do TOD no ambiente escolar e familiar.
    \item Identificar as principais dificuldades enfrentadas por pais e professores na comunicação sobre o comportamento das crianças.
    \item Propor funcionalidades tecnológicas que favoreçam o registro, acompanhamento e compartilhamento de informações relevantes.
    \item Avaliar a eficácia do sistema de informação como ferramenta de apoio à inclusão e ao desenvolvimento infantil.
\end{itemize}

\vspace{0.3cm}
\section{Estrutura do Projeto}


O projeto está organizado em capítulos que apresentam, de forma clara e sequencial, os principais aspectos do estudo. Inicialmente, o capítulo 1 contextualiza o tema, define o problema de pesquisa, os objetivos, a justificativa e a relevância do trabalho. Em seguida, o Capítulo 2 dedica-se à pesquisa de trabalhos correlatos, analisando sistemas similares para estabelecer um comparativo com a proposta desenvolvida. No capítulo 3, a  metodologia detalha os métodos adotados para a pesquisa, o desenvolvimento do sistema e os critérios para avaliação. O Capítulo 4 mostra a proposta de desenvolvimento do sistema e descreve as funcionalidades, a arquitetura e a implementação da ferramenta proposta, apresentando o levantamento de requisitos, os protótipos da interface e o cronograma de atividades que guiará a implementação e a avaliação da plataforma no próximo semestre. Por fim, o Capítulo 5 apresenta as considerações finais da pesquisa, sintetizando os resultados obtidos com a fase de ideação e prototipagem, além de apontar os direcionamentos e os trabalhos futuros para a continuidade do projeto.

\vspace{0.3cm}
\section{Classificação da Pesquisa}

A pesquisa classifica-se como aplicada, pois busca a solução de um problema prático relacionado ao manejo comportamental de crianças com TDAH e TOD por meio do desenvolvimento de um sistema de informação. Além disso, caracteriza-se como uma pesquisa qualitativa, uma vez que envolve a compreensão dos comportamentos, das interações e das percepções dos pais e professores, bem como a avaliação da eficácia do sistema proposto no contexto educacional. Essa abordagem permite uma análise aprofundada dos fenômenos envolvidos, valorizando a experiência dos participantes e a contextualização dos dados obtidos.
